\documentclass[11pt,a4paper]{article}

\usepackage[utf8]{inputenc}
\usepackage[T1]{fontenc}
\usepackage{amsmath,amssymb,amsfonts}
\usepackage{graphicx}
\usepackage{xcolor}
\usepackage[numbers,sort&compress]{natbib}
\usepackage[margin=2.2cm]{geometry}
\usepackage{booktabs}
\usepackage{microtype}
\usepackage{url}
\usepackage{hyperref}
\usepackage{doi}
\usepackage{times}

\hypersetup{
    colorlinks=true,
    linkcolor=blue!70!black,
    citecolor=blue!70!black,
    urlcolor=blue!70!black,
    pdftitle={Opinion–conviction decoupling explains late-campaign electoral swings and how socio-technical counter-strategies neutralize them},
    pdfauthor={César Rennó-Costa}
}

\title{\textbf{\Large Opinion–conviction decoupling explains late-campaign electoral swings and how socio-technical counter-strategies neutralize them}}

\author{
    \textbf{César Rennó-Costa}$^{1,*}$ \\[0.8em]
    $^{1}$Digital Metropolis Institute, Federal University of Rio Grande do Norte, Natal, RN, Brazil \\
    $^{*}$Corresponding author. Email: \texttt{cesar@imd.ufrn.br}
}

\date{}

\begin{document}

\maketitle

\begin{abstract}
Voters' core political convictions exhibit strong temporal stability, yet democratic elections can experience abrupt late-campaign volatility that confounds pre-election polling. While single-round polling discrepancies have occurred across major democracies, experiencing consecutive late-campaign reversals in opposite directions across sequential rounds is an empirical anomaly. Here I analyze this two-round divergence using an empirical digital twin of the 2026 Brazilian presidential election, where an eve-of-election conservative surge was followed by a runoff reversal. Coupling a census-calibrated urban multiplex network, broadcast media, and algorithmic platforms, the model reproduces longitudinal survey trajectories and official ballot totals ($R^2 = 0.9955$, root-mean-square error $0.85$ percentage points). I show that fast-opinion elasticity temporarily decouples expressed voting choices from stable convictions. Amplified by asymmetric candidate aversion ($3.5\times$ faster accumulation than decay), this decoupling enables dark-social cascades to overturn leads during polling blackouts, while coordinated socio-technical counter-strategies can rapidly neutralize these swings.
\end{abstract}

\vspace{0.8em}
\noindent\textit{Keywords:} Computational social science; Agent-based digital twin; Opinion dynamics; Algorithmic polarization; Political disinformation; Democratic resilience.

\vspace{1.5em}

A central tenet of political behavior is that citizens' core ideological convictions exhibit high temporal inertia, sustained by durable socioeconomic identities, generational socialization, and stable partisan attachments \citep{converse1964nature, alwin1991aging, zaller1992nature}. Over multi-year horizons, deep-seated political values change slowly, acting as psychological anchors that buffer democratic electorates against transient informational shocks \citep{kiley2020measuring, howe2017attitude}. Yet high-profile democratic contests have repeatedly drawn widespread attention for abrupt, late-breaking swings that defy pre-election survey projections despite the documented stability of underlying ideological convictions \citep{tucker2018social, lorenz2019accelerating, jennings2018election}. Notable precedents include late-stage shifts observed in the 2016 United Kingdom European Union membership referendum and the 2016 United States presidential election \citep{kennedy2018evaluation, jennings2018election}, as well as Brazil's 2018 presidential surge \citep{nicolau2020brasil} and the 2022 Brazilian first round, where incumbent Jair Bolsonaro outperformed eve-of-election polling estimates by over five percentage points \citep{tse2022resultados}. However, while these historical contests involved single late-campaign polling misses in one-round or decisive general elections, experiencing two consecutive late-campaign reversals in opposite directions across sequential rounds of the same presidential contest represents a rare and striking empirical anomaly.

Deciphering the mechanisms of late-campaign volatility requires moving beyond single-channel models of public opinion. Real-world political deliberation unfolds across three distinct, dynamically coupled communication substrates \citep{boccaletti2014structure, kivela2014multilayer}. First, face-to-face interpersonal interactions across everyday urban environments (households, workplaces, schools, religious congregations, friendship circles, and transit spaces) provide localized social reinforcement and cross-cutting deliberative friction \citep{granovetter1973strength, mcpherson2001birds}. Adapted from census-calibrated multi-layer urban agent-based models developed in computational epidemiology \citep{lopes2022measuring}, realistic sociopolitical digital twins must capture these structured contact environments to represent interpersonal contagion across city scales. Second, traditional mass media institutions (broadcast television, radio, and investigative print journalism) synchronize collective attention and establish professional editorial gatekeeping \citep{mccombs1972agenda, gentzkow2006media}. Third, modern algorithmic social platforms and dark-social encrypted messaging networks curate personal discovery feeds via engagement-maximizing recommendation engines while enabling frictionless peer-to-peer forwarding \citep{cinelli2021echo, vosoughi2018spread, cesarino2020posverdade}. 

A critical theoretical gap concerns how these three communicative substrates interact under acute, time-compressed informational offensives. Traditional campaign institutions rely on statutory broadcast airtime, legacy press conferences, and judicial fact-checking petitions \citep{gentzkow2006media, chadwick2017hybrid}. However, empirical evidence indicates that these defensive mechanisms suffer from a wide separation of operational timescales \citep{downs1972up, harder2017intermedia}: while institutional newsrooms and regulatory courts deliberate over days or weeks, algorithmic social networks propagate moral-emotional cascades in minutes \citep{brady2017emotion, vosoughi2018spread}. Furthermore, contemporary feed recommendation engines optimize engagement symmetrically for both positive affirmation and negative outrage \citep{crockett2017moral, bail2018exposure}. Consequently, defensive debunking and quote-dunking paradoxically supply the very engagement signals required by curation algorithms to amplify opponent disinformation into unaligned discovery feeds.

The 2026 Brazilian presidential election provides a uniquely revealing quasi-natural experiment to investigate these dynamics. Unlike previous contests characterized by a single polling miss, Brazil in 2026 exhibited a genuine double polling anomaly: tracking surveys failed to capture an acute election-eve conservative surge in the first round and subsequently missed a progressive counter-surge in the second round. Pitting incumbent president Luiz Inácio Lula da Silva against Senator Flávio Bolsonaro, the contest took place in a mobile-first democracy where statutory television advertising and digital networks operated across two distinct regulatory regimes. Throughout the first round, the incumbent coalition maintained a broadcast-centric strategy backed by a $+27.3\%$ advantage in statutory television airtime. Despite steady polling leads over fifty days, the incumbent was overtaken on election eve by an acute digital shock that lifted Flávio Bolsonaro into first place with $47.03\%$ of valid votes, confounding eve-of-election survey projections. In the ensuing runoff, statutory broadcast airtime was equalized to strict 50/50 parity by law, removing broadcast advantages and isolating digital campaign strategy as the primary differentiator. Faced with an initial deficit that tracking polls projected would solidify into defeat, the progressive campaign executed a strategic pivot: it adopted the counter-tactical doctrine of \emph{Janonismo Cultural}, articulated by Federal Deputy André Janones \citep{janones2023janonismo}, deploying socio-technical interventions designed to starve rival outrage, concentrate viral bandwidth, and bypass open feeds through dark-social messaging.

Here, I develop a high-resolution computational digital twin of the 2026 Brazilian electorate to systematically investigate the microscopic mechanisms of late-campaign electoral shocks and evaluate the efficacy of socio-technical counter-tactics. Calibrated to national census demography, detailed tracking surveys, and empirical platform telemetry, the agent-based model couples a five-layer multiplex interpersonal network, an institutional mass media ecosystem, and four algorithmic social platforms across a four-dimensional concept hypercube. By formalizing a two-timescale separation between fast-moving expressed opinions $\mathbf{x}_i(t)$ and slow-crystallizing convictions $\mathbf{s}_i(t)$, this framework captures how transient digital perturbations decouple public opinion from foundational beliefs. The same fast-opinion elasticity $\mathbf{x}_i(t)$ that leaves electorates vulnerable to acute late-campaign shocks also explains why coordinated socio-technical counter-strategies can rapidly neutralize those swings within a short runoff window without needing to reshape slow-moving deep convictions $\mathbf{s}_i(t)$. Precisely because 2026 Brazil exhibited both an initial late-stage conservative surge in Round~1 and a subsequent progressive reversal in Round~2, the contest provides an empirical setting to evaluate how the same opinion-elasticity mechanisms can both destabilize and restore an electoral majority. Through large-scale Monte Carlo ensembles ($N = 100{,}000$ citizens across 100 independent realizations), I isolate the psychological and topological drivers that enable late-stage disinformation blitzes to overturn established majorities during polling blackouts, and evaluate how structured algorithmic counter-offensives can mitigate late-stage erosion to restore electoral resilience.

\section*{Results}

\subsection*{Empirical dynamics of the 2026 presidential election}

The 2026 Brazilian presidential contest provides a detailed longitudinal record of electoral competition under divergent communicative regimes (Fig.~\ref{fig:fig1_election_data}). Under Brazilian electoral law (Lei n\textordmasculine~9.504/1997), presidential elections employ a two-round majoritarian system with strict statutory regulations governing campaign finance and broadcast airtime (Supplementary Note~4). In the first round, incumbent president Luiz Inácio Lula da Silva assembled a broad seven-party coalition securing 5m~31s per daily broadcast block ($44.1\%$ of total airtime, 433 thirty-second insertions), compared to 4m~20s ($34.7\%$, 339 insertions) for Senator Flávio Bolsonaro, generating a $+27.3\%$ broadcast endowment advantage for the incumbent (Fig.~\ref{fig:fig1_election_data}d). Financially, while both campaigns adhered to the statutory first-round spending ceiling of $\mathrm{R}\$$~88.9 million (Fig.~\ref{fig:fig1_election_data}e), they diverged sharply in their revenue models (Fig.~\ref{fig:fig1_election_data}f): Lula's campaign was $96.3\%$ funded by public transfers from the Special Campaign Financing Fund (FEFC; $\mathrm{R}\$$~123.8M of $\mathrm{R}\$$~128.6M declared), whereas Flávio's campaign deployed a hybrid model raising an unprecedented $\mathrm{R}\$$~43.2M in physical individual donations and micro-crowdfunding, establishing a $9.0\times$ private capital advantage over Lula ($\mathrm{R}\$$~4.8M).

Throughout the 50-day first-round campaign, voter intentions remained stable (Fig.~\ref{fig:fig1_election_data}b): longitudinal tracking surveys by Datafolha placed Lula ahead in a stable band of $41.8\%$--$45.0\%$ of valid votes against $37.5\%$--$41.0\%$ for Flávio, with third-way contenders holding $14.0\%$--$22.0\%$. Concurrently, candidate rejection rates remained mutual and stable between $45\%$ and $47\%$ (Fig.~\ref{fig:fig1_election_data}c). However, the first-round ballot on October 4 produced an abrupt empirical shock (Fig.~\ref{fig:fig1_election_data}a): Flávio Bolsonaro finished in first place with $47.03\%$ of valid votes ($55.7$ million votes), surpassing Lula ($45.16\%$, $53.5$ million votes) as third-party support collapsed to $7.81\%$ ($9.3$ million votes). Concurrently, voter abstention and blank ballots surged to a historic high of $24.70\%$ ($20.90\%$ abstention, $4.80\%$ blank/null; Fig.~\ref{fig:fig1_election_data}a). On election night, conservative digital channels mobilized with high intensity, generating $3.90$ million daily interactions compared to Lula's $0.17$ million, establishing a $22.9\times$ engagement advantage (Fig.~\ref{fig:fig1_election_data}h).

\begin{figure*}[t!]
\centering
\includegraphics[width=\textwidth]{figures/figure0_jan_electiondata.pdf}
\caption{\textbf{Asymmetric campaign endowments and late-stage digital surges drove a two-round electoral reversal in Brazil in 2026.} 
\textbf{a}, Official Superior Electoral Court (TSE) election results for Round~1 (October 4) and Round~2 (October 25), alongside election-day turnout mobilization that compressed non-voters by $-2.91$ percentage points. 
\textbf{b}, Longitudinal tracking of valid vote intentions across six Round~1 and four Round~2 Datafolha tracking surveys, with gold markers indicating official ballot tallies. 
\textbf{c}, Candidate rejection dynamics (Datafolha: percentage stating they would not vote for the candidate under any circumstance), highlighting the second-round reversal that drove Flávio's rejection to $46.0\%$. 
\textbf{d}, Statutory free broadcast airtime (HGPE) per daily block on television and radio, shifting from a $+27.3\%$ coalition advantage for Lula in Round~1 to strict legal parity (5m~00s each) in Round~2. 
\textbf{e}, Macro public party financing (FEFC; total national pool $\mathrm{R}\$$~4.96 billion) and statutory TSE presidential spending ceilings ($\mathrm{R}\$$~88.9M in Round~1, $\mathrm{R}\$$~133.4M cumulative in Round~2). 
\textbf{f}, Presidential ticket funding profiles, comparing Lula's state-financed budget ($96.3\%$ FEFC) with Flávio's hybrid apparatus that secured a $9.0\times$ private donation advantage. 
\textbf{g}, Quaest Digital Popularity Index (IPD; 0--100 scale, 175 metrics) across three milestone waves, illustrating a $+16.0$-point net turnaround swing for Lula. 
\textbf{h}, Daily social interaction volumes (millions of interactions per day, Bites data) during the Round~1 shock, televised debate, and the 13-day runoff window. 
\textbf{i}, Follower ecosystem and bench architecture, reporting candidate direct accounts across five platforms, official party handles, closed WhatsApp groups (Exército do Lula 10.0M vs.\ Tropa do Flávio 2.0M), and the arithmetic mean of the top-5 legislative surrogates per coalition.}
\label{fig:fig1_election_data}
\end{figure*}

This outcome prompted an immediate political realignment. Initial post-election polls projected a clear runoff lead for Flávio ($52.0\%$ vs.\ $48.0\%$ on October 8; Fig.~\ref{fig:fig1_election_data}b), reinforced by a post-victory momentum that reduced his rejection to $43.0\%$ while elevating Lula's to $47.0\%$ (Fig.~\ref{fig:fig1_election_data}c), mirrored by a 9.3-point conservative lead in the Quaest Digital Popularity Index (80.0 vs.\ 70.7; Fig.~\ref{fig:fig1_election_data}g). Under statutory runoff rules, television and radio airtime was equalized to strict 5m~00s parity (25 spots/day each; Fig.~\ref{fig:fig1_election_data}d), eliminating broadcast advantages and shifting the decisive contest to digital networks. Although Flávio maintained an elite surrogate bench strength averaging $17.4$ million followers per ally (led by Nikolas Ferreira with 52.0M, versus $7.1$ million for PT allies; Fig.~\ref{fig:fig1_election_data}i), Lula deployed a personal audience of $39.0$ million followers and a synchronized closed messaging infrastructure of $10.0$ million participants (Exército do Lula; Fig.~\ref{fig:fig1_election_data}i).

During the runoff, the progressive campaign adopted an active digital counter-offensive. Following the October 16 televised debate, progressive channels generated $25.5$ million interactions against $5.0$ million for the opponent ($5.1\times$ ratio; Fig.~\ref{fig:fig1_election_data}h), maintaining a sustained $3.5\times$ daily interaction advantage throughout the runoff ($8.52$M vs.\ $2.41$M/day). This sustained engagement steadily eroded Flávio's momentum: by October 22, Flávio's rejection rose to $46.0\%$ while polling intentions reached a dead heat ($50.0\%$ vs.\ $50.0\%$; Fig.~\ref{fig:fig1_election_data}b,c), accompanied by a $+16.0$-point net swing in the Quaest IPD (Lula $+6.7$, $83.3$ vs.\ $76.6$; Fig.~\ref{fig:fig1_election_data}g). On election day (October 25), assisted by municipal free public transit decrees that lowered economic voting frictions, voter non-participation fell by $2.91$ percentage points (from $24.70\%$ to $21.79\%$; Fig.~\ref{fig:fig1_election_data}a). Lula secured victory with $51.79\%$ of valid votes ($61.8$M vs.\ $57.5$M), representing the first time in modern Brazilian democratic history that a presidential candidate trailing in the first round won the executive in the runoff.

\subsection*{Baseline relaxation physics of urban attitude networks}

To establish the baseline stability of public attitudes prior to media exposure or algorithmic amplification, I first quantify the relaxation physics of the urban interpersonal network (Fig.~\ref{fig:fig2_interpersonal_network}). Calibrated to IBGE 2022 National Population Census demography \citep{ibge2022censo}, the model situates citizens in a four-dimensional concept hypercube $\Omega = [-1, 1]^4$ \citep{converse1964nature}. Individual behavioral parameters, cognitive plasticity $\alpha_i$, bounded confidence threshold $\epsilon_i$, attitude uncertainty $u_i$, and confirmation bias $b_i$, are conditioned on empirical age and socioeconomic status (SES; Fig.~\ref{fig:fig2_interpersonal_network}a--c). In agreement with the aging-stability hypothesis \citep{alwin1991aging, kiley2020measuring}, cognitive plasticity $\alpha(a)$ undergoes monotonic decay with age (Fig.~\ref{fig:fig2_interpersonal_network}a), while older and lower-SES cohorts exhibit narrower confidence bounds $\epsilon_i$ and elevated confirmation bias $b_i$ (Supplementary Note~1).

\begin{figure*}[t!]
\centering
\includegraphics[width=\textwidth]{figures/figure1_jan_nonmedianetwork.pdf}
\caption{\textbf{Decentralized deliberation across urban multiplex networks converges toward a scale-invariant centrist equilibrium while buffering core convictions against acute shocks.} 
\textbf{a}, Empirical age pyramid from the 2022 IBGE Census and age-dependent cognitive plasticity $\alpha(a)$. 
\textbf{b}, Urban socioeconomic status (SES) distribution fitted to a Beta(2, 3) density. 
\textbf{c}, Cohort distributions of individual behavioral traits ($\alpha, \epsilon, u, b$). 
\textbf{d}, Five-layer multiplex network topology $\mathcal{M} = \{\mathcal{L}_1, \dots, \mathcal{L}_5\}$ spanning household, workplace/school, religious congregations, friendship, and transit weak ties. 
\textbf{e}, Inter-cohort age contact matrix $M_{ab}$ across the six adult age brackets. 
\textbf{f}, Age homophily $h_{\mathrm{age}}$ and ideological assortativity $r_{\mathrm{ideol}}$ across layers $\mathcal{L}_1$--$\mathcal{L}_5$. 
\textbf{g}, Two-dimensional projection of initial convictions $\mathbf{s}_i(0)$ and converged convictions $\mathbf{s}_i(T)$ alongside exogenous facts $\mathbf{f}_k$. 
\textbf{h}, Temporal evolution of pairwise cross-issue alignment $\langle |\rho_{kl}(t)| \rangle$ across city sizes ($N = 10\mathrm{k}$ to $10\mathrm{M}$). 
\textbf{i}, Asymptotic alignment ($\max |\rho|$, $\langle |\rho| \rangle$) and effective ideological dimensionality $\mathrm{PR}(\Sigma)$ versus urban population scale. 
\textbf{j}, Stationary marginal probability densities $p(s_d)$ across ideological dimensions $D_1$--$D_4$. 
\textbf{l}, Two-timescale relaxation of mean opinion $\bar{\mathbf{x}}(t)$ and conviction $\bar{\mathbf{s}}(t)$ under a transient exogenous shock. 
\textbf{m}, Age-stratified opinion--conviction tether stretching $\|\mathbf{x}_i(t) - \mathbf{s}_i(t)\|$ during the shock window. 
\textbf{n}, Attitude space snapshot at peak shock excursion ($t_{\mathrm{shock}}$), contrasting displaced opinions $\mathbf{x}_i$ against anchored convictions $\mathbf{s}_i$.}
\label{fig:fig2_interpersonal_network}
\end{figure*}

Interpersonal interactions unfold over a five-layer multiplex topology $\mathcal{M} = \{\mathcal{L}_1, \dots, \mathcal{L}_5\}$ \citep{boccaletti2014structure, kivela2014multilayer} capturing household cohabitation ($\mathcal{L}_1$, mean size $2.79$, weight $w_1 = 0.36$), workplace/school cohorts ($\mathcal{L}_2$, $\langle k^{[2]} \rangle = 10$, $w_2 = 0.26$), religious and community congregations ($\mathcal{L}_3$, $\langle k^{[3]} \rangle = 8$, $w_3 = 0.14$), small-world friendship circles ($\mathcal{L}_4$, $\langle k^{[4]} \rangle = 6$, $w_4 = 0.14$), and casual transit ties ($\mathcal{L}_5$, $\langle k^{[5]} \rangle = 4$, $w_5 = 0.10$) (Fig.~\ref{fig:fig2_interpersonal_network}d). The empirical age contact matrix $M_{ab}$ balances generational stratification in workplaces against cross-cohort mixing in households and religious institutions (Fig.~\ref{fig:fig2_interpersonal_network}e). This architecture generates a structural trade-off between age homophily $h_{\mathrm{age}}$ and ideological assortativity $r_{\mathrm{ideol}}$ (Fig.~\ref{fig:fig2_interpersonal_network}f): while friendship ($\mathcal{L}_4$, $r_{\mathrm{ideol}} = 0.38$) and religious groups ($\mathcal{L}_3$, $r_{\mathrm{ideol}} = 0.46$) promote ideological sorting, transit encounters ($\mathcal{L}_5$, $r_{\mathrm{ideol}} \approx 0.02$) and domestic ties ($\mathcal{L}_1$, $r_{\mathrm{ideol}} = 0.24$) act as cross-cutting bridges that inhibit macroscopic topological fragmentation.

In the absence of external media distortion, decentralized deliberation over this multiplex substrate systematically dampens ideological extremity (Fig.~\ref{fig:fig2_interpersonal_network}g). Although initial convictions $\mathbf{s}_i(0)$ are widely dispersed across $[-1, 1]^4$, bounded-confidence exchanges pull terminal convictions $\mathbf{s}_i(T)$ toward a cohesive centrist consensus around the origin $\mathbf{0}$ \citep{castellano2009statistical}. Scaling analysis across four orders of magnitude in urban population ($N = 10^4$ to $10^7$) indicates that this macroscopic convergence is scale-invariant (Fig.~\ref{fig:fig2_interpersonal_network}h--j). For populations $N \ge 100\mathrm{k}$, finite-size stochasticity vanishes into a thermodynamic limit where terminal issue correlation ($\max |\rho| \approx 0.08, \langle |\rho| \rangle \approx 0.03$) and effective belief dimensionality ($\mathrm{PR}(\Sigma) \approx 3.75$ out of $D=4$) stabilize at identical asymptotic values (Fig.~\ref{fig:fig2_interpersonal_network}i), while stationary marginal densities $p(s_d)$ remain unimodal and Gaussian-like across all dimensions (Fig.~\ref{fig:fig2_interpersonal_network}j).

This formulation resolves the dynamical coupling between fast expressed opinions $\mathbf{x}_i(t)$ and slow internalized convictions $\mathbf{s}_i(t)$ \citep{zaller1992nature, howe2017attitude} (Fig.~\ref{fig:fig2_interpersonal_network}l--n). Under an acute exogenous factual shock, mean opinion $\bar{\mathbf{x}}(t)$ responds rapidly ($\tau_x \approx 0.15$), whereas mean conviction $\bar{\mathbf{s}}(t)$ exhibits strong structural inertia ($\tau_s / \tau_x \approx 11.3$), shifting only marginally before pulling opinions back toward baseline during post-shock relaxation (Fig.~\ref{fig:fig2_interpersonal_network}l). This separation of timescales induces an elastic tether displacement $\|\mathbf{x}_i(t) - \mathbf{s}_i(t)\|$ that varies across demographics (Fig.~\ref{fig:fig2_interpersonal_network}m): younger citizens experience substantial opinion displacement due to higher cognitive plasticity, while older citizens remain tightly tethered to crystallized convictions \citep{alwin1991aging}. A cross-sectional snapshot of the belief hypercube at peak shock excursion ($t_{\mathrm{shock}}$; Fig.~\ref{fig:fig2_interpersonal_network}n) illustrates this instantaneous decoupling between displaced opinions and stationary convictions, providing a mechanistic explanation for how acute campaign events can trigger large short-term shifts in voting intention without immediately destabilizing core ideological convictions.

\subsection*{Media market structure and editorial gatekeeping}

Introducing the second communicative substrate, institutional mass media channels structure collective public attention via centralized broadcast scheduling and editorial gatekeeping \citep{mccombs1972agenda} (Fig.~\ref{fig:fig3_traditional_media}). I parameterize five distinct modalities: Broadcast Television ($v = 0$), Radio ($v = 1$), Print and Legacy News Portals ($v = 2$), Outdoor Signage ($v = 3$), and Campaign Rallies ($v = 4$) (Fig.~\ref{fig:fig3_traditional_media}a). Modalities differ in their characteristic attention decay rate $\delta_v$ and institutional credibility $\gamma_v$ (Supplementary Note~2): Written Press combines high institutional credibility ($\gamma = 0.82$) with slow editorial decay ($\delta = 0.18$), Broadcast Television provides mass reach with intermediate persistence ($\gamma = 0.75, \delta = 0.30$), and Campaign Rallies deliver high emotional mobilization pulses ($\gamma = 0.80$) subject to rapid post-event decay ($\delta = 0.70$). Audience penetration curves $P_v(\mathrm{age})$ calibrated to Kantar IBOPE 2023 surveys \citep{kantar2023inside} indicate that linear television viewing increases monotonically with age, dominating news consumption among senior citizens, whereas Radio and Outdoor signage reach daily commuters (Fig.~\ref{fig:fig3_traditional_media}b).

\begin{figure*}[t!]
\centering
\includegraphics[width=\textwidth]{figures/figure2_jan_medianetwork.pdf}
\caption{\textbf{Editorial tendentiousness and market concentration nonlinearly govern collective factual accuracy and societal polarization.} 
\textbf{a}, Attention decay rates $\delta_v$ and institutional credibility $\gamma_v$ across five classical media modalities. 
\textbf{b}, Age-dependent demographic reach curves $P_v(\mathrm{age})$ calibrated to empirical audience surveys. 
\textbf{c}, Editorial projection of exogenous facts $\mathbf{f}_k$ toward media group centers $\mathbf{m}_g$ under factual ($\tau = 0.10$) and tendentious ($\tau = 0.90$) regimes. 
\textbf{d}, Population mean attitude trajectory $\bar{x}_1(t)$ following an acute factual shock ($\mathbf{f}^* = +0.80$ at $t \ge 25$) across five media market configurations. 
\textbf{e}, Ideological shift (top) and Esteban--Ray societal polarization $P_{\mathrm{ER}}$ (bottom) across market concentration regimes ($M \in \{0, 1, 2, 3, 4, \text{Many}\}$). 
\textbf{f}, Timescale separation between transient broadcast pulses ($\tau_{\mathrm{media}} \approx 4.8$~steps) and ongoing interpersonal network deliberation ($\tau_{\mathrm{network}} \approx 19.0$~steps).}
\label{fig:fig3_traditional_media}
\end{figure*}

Media conglomerates project factual notices toward editorial ideological centers $\mathbf{m}_g \in \Omega$ through their tendentiousness parameter $\tau_g \in [0, 1]$ \citep{mullainathan2005market, gentzkow2006media} (Methods). In attitude space (Fig.~\ref{fig:fig3_traditional_media}c), low tendentiousness ($\tau = 0.10$) confines published notices to a tight envelope around true factual coordinates $\mathbf{f}_k$. Conversely, under high tendentiousness ($\tau = 0.90$), editorial gatekeeping displaces published coverage toward the conglomerate's ideological anchor $\mathbf{m}_g$ \citep{puglisi2015newspaper}, substituting factual reality with partisan narratives.

Market structure nonlinearly governs collective responsiveness following an acute factual shock ($\mathbf{f}^* = +0.80$ introduced at $t \ge 25$; Fig.~\ref{fig:fig3_traditional_media}d). In the absence of broadcast media ($M=0$), information diffuses sluggishly through word-of-mouth alone, producing negligible collective adjustment ($\bar{x}_1 \approx 0.00 \to 0.03$). Factual broadcasting ($\tau = 0.10$) rapidly synchronizes public opinion toward empirical reality ($\bar{x}_1 \to +0.50$). Conversely, a biased media monopoly ($M=1, \tau = 0.90$ anchored at $\mathbf{m}_1 = -0.60$) actively drives public opinion in the opposite direction ($\bar{x}_1 \to -0.24$), inducing larger systemic error than the absence of media. Across market configurations ($M \in \{0, 1, 2, 3, 4, \text{Many}\}$; Fig.~\ref{fig:fig3_traditional_media}e), increasing outlet plurality dilutes monopoly capture (Fig.~\ref{fig:fig3_traditional_media}e, top). However, societal division, measured via the Esteban--Ray polarization index $P_{\mathrm{ER}}$ \citep{esteban1994measurement}, peaks sharply under a tendentious duopoly ($M=2$; Fig.~\ref{fig:fig3_traditional_media}e, bottom), which bifurcates citizens into polarized information silos. Only pluralistic markets ($M \ge 4$) disperse partisan sorting and restore low societal polarization comparable to the baseline urban network.

Coupling the broadcast layer to the interpersonal network reveals a fundamental separation of operational timescales (Fig.~\ref{fig:fig3_traditional_media}f). Broadcast notices display high instantaneous reach but rapid temporal decay ($\tau_{\mathrm{media}} \approx 4.8$~steps) \citep{downs1972up}. In contrast, interpersonal deliberation across the multiplex network operates on a substantially slower relaxation timescale ($\tau_{\mathrm{network}} \approx 19.0$~steps) \citep{friedkin1990social}. Consequently, when broadcast channels introduce a salient campaign event, its societal impact is sustained long after the initial broadcast window closes, as citizens continue to process the information across household and peer ties.

\subsection*{Algorithmic outrage amplification in digital platforms}

Completing the communication architecture, the digital layer introduces algorithmic social platforms and closed encrypted messaging networks \citep{cinelli2021echo, tucker2018social} (Fig.~\ref{fig:fig4_social_media}). Calibrated to national surveys \citep{reuters2023digital}, age access profiles $P_p(\mathrm{age})$ define four digital modalities (Fig.~\ref{fig:fig4_social_media}a): short-form visual discovery feeds (Platform~0: Instagram and TikTok), capturing over $88\%$ of youth cohorts; video commentary and microblogging (Platform~1: YouTube and X), engaging politically attentive citizens; reciprocal social feeds (Platform~2: Facebook), concentrating among mature adults; and closed encrypted messaging (Platform~3: WhatsApp and Telegram), operating as a universal communications backbone ($>80\%$ penetration across all adult brackets) structured into private micro-groups ($K \le 256$~members).

\begin{figure*}[t!]
\centering
\includegraphics[width=0.82\textwidth]{figures/figure3_jan_socialmedianetwork.pdf}
\caption{\textbf{Engagement-maximizing recommendation algorithms amplify polarizing content and sustain multi-stage cascades into encrypted messaging networks.} 
\textbf{a}, Age-dependent daily access probability $P_p(\mathrm{age})$ across four digital platform modalities and broadcast television. 
\textbf{b}, Heavy-tailed follower distribution $P(k) \sim k^{-\alpha}$ ($\alpha \approx 1.15$, Gini coefficient $0.73$) and emergent influencer hierarchy. 
\textbf{c}, Algorithmic relevance score $R_m$ mapped across positive-valence (admiration) and negative-valence (outrage) engagement coordinates, highlighting maximal amplification along the diagonal of controversy. 
\textbf{d}, Three-stage cross-platform meme diffusion lifecycle from visual discovery feeds (Stage~1) and public microblogging (Stage~2) to closed encrypted messaging cascades (Stage~3). 
\textbf{e}, Hybrid media agenda-setting dynamics following an authoritative investigative news release at $t=0$, contrasting legacy broadcast turnaround ($\tau_{\mathrm{turnaround}}$) against rapid digital response ($\tau_{\mathrm{social}}$). 
\textbf{f}, Long-horizon evolution of Esteban--Ray societal polarization $P_{\mathrm{ER}}(t)$ across creator tendentiousness regimes ($\tau_{\mathrm{inf}} \in \{0.10, \text{Empirical}, 0.90\}$).}
\label{fig:fig4_social_media}
\end{figure*}

Within open digital feeds, follower accumulation produces extreme structural inequality (Fig.~\ref{fig:fig4_social_media}b). Preferential attachment yields a scale-free in-degree distribution $P(k) \sim k^{-\alpha}$ with exponent $\alpha \approx 1.15$ and an empirical Gini coefficient of $0.73$ (Fig.~\ref{fig:fig4_social_media}b). Under this topology, the top $5\%$ of elite creators capture $54.5\%$ of direct audience reach, whereas the bottom $50\%$ command only $8.4\%$. This hyper-concentration converts decentralized networks into a star-like broadcasting system: when an elite influencer publishes a narrative, their message bypasses local interpersonal clustering to reach hundreds of thousands of individual algorithmic feeds simultaneously.

Feed recommendation algorithms rank content by optimizing for aggregate user engagement rather than factual accuracy \citep{vosoughi2018spread, crockett2017moral}. Computing relevance scores symmetrically across positive and negative interactions creates a structural amplification along the diagonal of controversy (Fig.~\ref{fig:fig4_social_media}c). Posts provoking both in-group admiration and out-group outrage receive maximal algorithmic distribution, rewarding polarizing moral-emotional framing \citep{brady2017emotion, crockett2017moral}.

This architecture generates a three-stage cross-platform cascade lifecycle (Fig.~\ref{fig:fig4_social_media}d; Supplementary Note~3). Diffusion initiates in Stage~1 on visual discovery feeds (Instagram/TikTok), surfacing novel memes to high-plasticity youth within $10$~hours. In Stage~2, narratives migrate to commentary platforms (YouTube/X), where quote-dunking triggers secondary viral waves between hours $15$ and $25$. In Stage~3, content penetrates closed messaging groups (WhatsApp/Telegram). Forwarded across trusted peer contacts ($\tau_{\text{WA}} = 0.85$, $2.8\times$ persuasion efficiency), dark-social cascades bypass platform moderation and persist long after open-feed attention has decayed. Coupling the digital layer to traditional newsrooms creates a hybrid agenda-setting feedback loop \citep{chadwick2017hybrid, harder2017intermedia} (Fig.~\ref{fig:fig4_social_media}e): while investigative journalism requires an editorial turnaround of $\tau_{\mathrm{turnaround}} \approx 1$--$2$~days, digital platforms respond within hours ($\tau_{\mathrm{social}} \approx 0.1$--$0.3$~days), generating social volume that compels broadcast media to maintain follow-up coverage. When influencers exploit this dynamic with high tendentiousness ($\tau_{\mathrm{inf}} = 0.90$), societal polarization climbs to an entrenched plateau of $P_{\mathrm{ER}} = 0.18$ (Fig.~\ref{fig:fig4_social_media}f).

\subsection*{Negativity asymmetry in late-campaign voting swings}

The convergence of multiplex interpersonal clustering, traditional media gatekeeping, and algorithmic outrage explains the dynamics of the first-round election (Fig.~\ref{fig:fig5_first_round}). Calibrating initial ideological identification to national survey benchmarks \citep{datafolha2024identificacao} establishes four distinct voting blocs (Fig.~\ref{fig:fig5_first_round}a): Lulistas ($\sim 40\%$, comprising $28\%$ convinced identifiers and $12\%$ non-convinced sympathizers), Bolsonaristas ($\sim 34\%$, $23\%$ convinced and $11\%$ sympathizers), Centrists ($\sim 19\%$), and unaligned non-voters ($\sim 7\%$). In concept space $\Omega$, Lula is positioned at Center-Left ($\mathbf{c}_L = [-0.25, -0.20, -0.35, -0.15]^\top$), Flávio Bolsonaro at Far-Right ($\mathbf{c}_F = [0.65, 0.70, 0.60, 0.65]^\top$), and Third-Way contenders at Center-Right ($\mathbf{c}_O = [0.20, 0.25, 0.10, 0.30]^\top$). Brazilian centrists are culturally conservative on moral and security axes ($d=1, 3$) yet supportive of state welfare and minimum wages ($d=0, 2$), making them a cross-pressured swing constituency \citep{almeida2019evangelicals}.

\begin{figure*}[p]
\centering
\includegraphics[width=0.92\textwidth,height=0.695\textheight,keepaspectratio]{figures/figure4_jan_election_first.pdf}
\caption{\textbf{Cognitive negativity asymmetry and a 48-hour dark-social blitz overturned an established polling lead in the first round.} 
\textbf{a}, Sociopolitical segmentation of the electorate into ideological identification blocs: Lulistas ($\sim 40\%$, decomposed into convict identifiers vs.\ non-convict sympathizers), Bolsonaristas ($\sim 34\%$, convict vs.\ non-convict), Centrists/Others ($\sim 19\%$), and unaligned/None ($\sim 7\%$), calibrated to national survey data. 
\textbf{b}, Longitudinal 50-day first-round campaign trajectories (Model Fit, solid curves vs.\ Neutral Model organic baseline, dashed curves) benchmarked against Datafolha tracking survey waves (diamonds $\pm 1\,\mathrm{SD}$) and official TSE ballot totals (circles on October 4: Flávio $47.0\%$ / $55.7\mathrm{M}$, Lula $45.2\%$ / $53.5\mathrm{M}$, Others $7.8\%$ / $9.3\mathrm{M}$), highlighting the August 23 debate shock ($*$) and the acute 48-hour pre-election blitz (shaded amber band). 
\textbf{c}, Coupled affective trajectory of population-mean candidate admiration (solid curves) and rejection (dashed curves) for Lula (red) and Flávio (blue), showing the acute late-campaign aversion spike triggered against Lula ($\lambda_R/\delta_R = 3.5\times$). 
\textbf{d}, Factorial attribution of Flávio's valid vote share shift ($\Delta V$) across the five Janones mechanisms: M1 Religious Priming ($+2.10\mathrm{M}$ votes), M2 Media Lean ($+1.42\mathrm{M}$), M3 Bot Astroturfing ($+1.60\mathrm{M}$), M4 Rejection News ($+1.84\mathrm{M}$), and M5 48-Hour WhatsApp Blitz ($+2.61\mathrm{M}$); error bars indicate the full simulation range $[\min, \max]$ across 100 independent Monte Carlo realizations ($N=100{,}000$). 
\textbf{e}, Multi-feature parametric dose-response showing Flávio vote share gain ($+\Delta V_F$, solid curves) and concurrent Lula share loss ($-\Delta V_L$, dashed curves) across normalized mechanism intensity ($0\text{--}100\%$ of maximum scale; physical bounds noted in axis annotation), with empirical operating points marked by circles (Flávio) and squares (Lula), referenced against the Status Quo Baseline ($\Delta V = 0$) and the $+1.8\%$ ($+2.13\mathrm{M}$ votes) iso-gain benchmark contour. 
\textbf{f}, Cumulative model stepwise build-up across stages (Neutral Baseline, +Church, +Media, +Bots, +Rejection, +Blitz), illustrating valid vote shares for Flávio Bolsonaro (blue) and Lula da Silva (red); error bars denote the full simulation range $[\min, \max]$ across 100 independent seeds ($N=100{,}000$), indicating that decisive lead reversal occurs exclusively upon the final 48-hour blitz.}
\label{fig:fig5_first_round}
\end{figure*}

Simulating the 50-day first round across 100 independent Monte Carlo realizations ($N = 100{,}000$ citizens) reproduces Datafolha survey waves and the official TSE outcome ($47.0\%$ Flávio, $45.2\%$ Lula, $7.8\%$ Others) with high statistical fidelity ($\mathrm{RMSE} = 0.85$ percentage points, $R^2 = 0.9955$; Fig.~\ref{fig:fig5_first_round}b). Under the neutral organic baseline (dashed curves), Lula preserves an uninterrupted lead throughout the campaign, concluding at $48.2\%$ ($57.1\mathrm{M}$ votes) versus $38.8\%$ for Flávio ($46.0\mathrm{M}$ votes). In the fitted model (solid curves), Lula leads securely through day 47 ($48.2\%$ vs $38.8\%$). During the final 48-hour sprint (shaded amber band), a synchronized multi-platform offensive produces a sharp vote shift of $+8.20$ percentage points ($+9.72\mathrm{M}$ votes), lifting Flávio ahead on election day.

The psychological engine powering this late-stage divergence is a cognitive negativity asymmetry governing candidate admiration and rejection (Fig.~\ref{fig:fig5_first_round}c; Methods). While admiration accumulates slowly and decays rapidly, candidate aversion accumulates rapidly and decays with high hysteresis. This kinetic asymmetry, where aversion accumulates $3.5$ times faster than it decays, implies that negative shocks inflict persistent psychological scars that resist correction. During the final 48 hours, targeted moral panics elevate Lula's rejection from $0.38$ to $0.50$ (Fig.~\ref{fig:fig5_first_round}c). Because voter utility penalizes candidate aversion more heavily than it rewards positive affinity (Methods), this rejection spike depresses swing-voter turnout and shifts cross-pressured centrists toward Flávio.

Factorial attribution isolates the specific contribution of each mechanism toward Flávio's $+8.20\%$ ($+9.72\mathrm{M}$ votes) surge (Fig.~\ref{fig:fig5_first_round}d): Religious Priming (M1, $+1.77\%$, $+2.10\mathrm{M}$ votes) anchors Evangelical voters in layer $\mathcal{L}_3$; Traditional Media Slant (M2, $+1.20\%$, $+1.42\mathrm{M}$ votes) reinforces fiscal conservatism; Bot Astroturfing (M3, $+1.35\%$, $+1.60\mathrm{M}$ votes) manufactures artificial outrage on open feeds; Asymmetric Rejection News (M4, $+1.55\%$, $+1.84\mathrm{M}$ votes) exploits cognitive negativity bias; and the 48-Hour WhatsApp Blitz (M5, $+2.20\%$, $+2.61\mathrm{M}$ votes) floods dark-social groups beyond journalistic verification horizons. Parametric dose-response curves reveal that the WhatsApp blitz exhibits the steepest marginal electoral return ($\mathrm{d}V_F/\mathrm{d}\theta \approx 0.034\%$ per percentage point; Fig.~\ref{fig:fig5_first_round}e). A cumulative build-up test reveals a threshold principle (Fig.~\ref{fig:fig5_first_round}f): Lula preserves his polling lead through stages 1 to 4 (+Church, +Media, +Bots, +Rejection); decisive lead reversal is triggered exclusively when the acute dark-social blitz is added in Stage 5, lifting Flávio to $47.0\%$ vs $45.2\%$.

\subsection*{Algorithmic counter-mobilization during the runoff}

The first-round diagnosis informed the design of the progressive counter-offensive in the second-round runoff (October 4--25, 2026; Fig.~\ref{fig:fig6_second_round}), which operationalized the strategic doctrine of \emph{Janonismo Cultural} \citep{janones2023janonismo}. In counterfactual simulations, if the progressive campaign had maintained its passive first-round posture relying on broadcast airtime and conventional fact-checking, the same-strategy counterfactual (dashed curves, Fig.~\ref{fig:fig6_second_round}a) indicates that Flávio would have consolidated centrist support to win decisively with $53.6\%$ ($63.9\mathrm{M}$ votes) against $46.4\%$ for Lula ($55.4\mathrm{M}$ votes).

\begin{figure*}[p]
\centering
\includegraphics[width=0.92\textwidth,height=0.725\textheight,keepaspectratio]{figures/figure5_jan_election_second.pdf}
\caption{\textbf{Coordinated socio-technical counter-tactics inverted candidate rejection trajectories and crossed the runoff majority threshold.} 
\textbf{a}, Longitudinal 21-day second-round calibration ($N = 100{,}000$ citizens, 100-seed ensemble $\pm 1\,\mathrm{SD}$ envelope) contrasting the Same-Strategy Counterfactual (Neutral Model, dashed curves: Flávio wins $53.6\%$ vs $46.4\%$) against the fitted Janonismo Cultural counter-offensive (Model Fit, solid curves: Lula wins $51.8\%$ / $61.8\mathrm{M}$ vs Flávio $48.2\%$ / $57.5\mathrm{M}$, jumping $+6.64\%$ from Round 1), incorporating the October 18 Band debate shock ($*$), the final 48-hour sprint (shaded amber band), Datafolha tracking surveys (diamonds $\pm 1\,\mathrm{SD}$), and official TSE runoff results (yellow-filled markers on October 25). 
\textbf{b}, Coupled affective trajectory illustrating the complete structural inversion of candidate rejection (dashed curves: Flávio $\bar{R}_F$ climbs $40\% \to 53\%$, while Lula $\bar{R}_L$ drops $50\% \to 41\%$) alongside admiration consolidation (solid curves: Lula $\bar{A}_L$ climbs $43\% \to 52\%$). 
\textbf{c}, Factorial attribution of Lula's $+6.64\%$ ($+7.93\mathrm{M}$ votes) runoff turnaround across the six Janonismo counter-tactics: T1 Algorithmic Asphyxiation ($+1.36\mathrm{M}$ votes), T2 Strict Monothematism ($+1.61\mathrm{M}$), T3 Material Loss-Aversion Framing ($+1.69\mathrm{M}$), T4 Molecular WhatsApp Synchronization ($+1.25\mathrm{M}$), T5 Mass Capital Rallies ($+1.05\mathrm{M}$), and T6 24-Hour Live Saturation ($+0.95\mathrm{M}$); error bars indicate the full simulation range $[\min, \max]$ across 100 independent Monte Carlo seeds ($N=100{,}000$). 
\textbf{d}, Multi-tactic parametric dose-response curves and iso-vote equivalence contours mapping normalized tactic intensity ($0\text{--}100\%$ of maximum scale; physical bounds noted in axis annotation) to Lula valid share gain ($+\Delta V_L$, $\%$), with calibrated empirical operating points and benchmark iso-gain levels marked at $+1.69\mathrm{M}$ ($+1.42\%$), $+1.36\mathrm{M}$ ($+1.14\%$), and $+1.05\mathrm{M}$ ($+0.88\%$). 
\textbf{e}, Stepwise cumulative model addition (R1 Base, Same Strategy, +T1, +T2, +T3, +T4, +T5, +T6) showing valid vote shares for Lula (red) and Flávio (blue); error bars denote the full simulation range across 100 independent seeds ($N=100{,}000$), illustrating the $50\%$ majority threshold crossover at Tactic 4 (WhatsApp Sync, $50.1\%$) and final victory consolidation at $51.8\%$.}
\label{fig:fig6_second_round}
\end{figure*}

Under the active counter-tactical suite, the campaign executed six structured interventions, formalizing the core concepts of \emph{Janonismo Cultural} into socio-technical mechanisms (Methods; Supplementary Note~5):
1. \emph{Outrage-Starvation Embargo} (T1: Algorithmic Asphyxiation, $+1.14\%$, $+1.36\mathrm{M}$ votes): Allied accounts enforced complete non-engagement toward rival provocations, depriving curation engines of controversy telemetry and collapsing opponent feed injection by $48\%$.
2. \emph{Monothematic Bandwidth Concentration} (T2: Strict Monothematism, $+1.35\%$, $+1.61\mathrm{M}$ votes): Allied communication bandwidth was concentrated onto a single negative theme per 24-hour cycle, exploiting superlinear returns in platform discovery functions to quadruple viral reach over dispersed messaging.
3. \emph{Loss-Aversion Framing} (T3: Material Loss Frame, $+1.42\%$, $+1.69\mathrm{M}$ votes): Replaced procedural defenses of democratic institutions with concrete threats to personal welfare (welfare rollbacks, minimum wage freezes), engaging Kahneman--Tversky loss weighting (a 2.25-fold aversion to losses relative to equivalent gains) to convert $52.4\%$ of undecided working-class swing voters.
4. \emph{Molecular Encrypted Cascades} (T4: WhatsApp Synchronization, $+1.05\%$, $+1.25\mathrm{M}$ votes): Coordinated a registered network of $10.0\mathrm{M}$ participants to dispatch curated multimedia assets into private family and neighborhood groups, leveraging high interpersonal trust in direct peer messaging to achieve a 2.8-fold conversion multiplier over open feeds.
5. \emph{Physical Turnout Locking} (T5: Mass Capital Rallies, $+0.88\%$, $+1.05\mathrm{M}$ votes): Organized mass public assemblies in 27 state capitals ($1.80\mathrm{M}$ attendees) combined with free transit decrees to eliminate travel friction, locking committed voter turnout and providing organic visual proof of popular mobilization.
6. \emph{Attention-Ceiling Saturation} (T6: 24-Hour Continuous Livestream, $+0.80\%$, $+0.95\mathrm{M}$ votes): Maintained an unbroken 24-hour broadcast stream during the final 48 hours capturing $65\%$ of public digital attention, imposing an attention ceiling that compressed rival blitz discovery below $20\%$.

Together, these interventions produced a complete inversion of candidate affective evaluations (Fig.~\ref{fig:fig6_second_round}b): Flávio's rejection rose from $\bar{R}_F = 0.40$ to $0.53$, while Lula's rejection dropped from $\bar{R}_L = 0.50$ to $0.41$, accompanied by a rise in Lula admiration ($\bar{A}_L = 0.43 \to 0.52$). The calibrated model reproduces the 21-day runoff tracking polls and official election results with high precision ($\mathrm{RMSE} = 0.56$ percentage points, $R^2 = 0.9128$; Fig.~\ref{fig:fig6_second_round}a). Parametric dose-response curves illustrate tactical fungibility across the intervention suite (Fig.~\ref{fig:fig6_second_round}d). A cumulative addition test reveals a decisive majority crossover (Fig.~\ref{fig:fig6_second_round}e): starting from the Round 1 baseline ($45.16\%$), tactics T1 ($46.30\%$), T2 ($47.65\%$), and T3 ($49.07\%$) steadily close the gap, while the $50\%$ majority threshold is crossed upon adding Tactic T4 (Molecular Encrypted Cascades, reaching $50.12\%$), with T5 ($51.00\%$) and T6 ($51.80\%$) consolidating the final victory.

\section*{Discussion}

The findings presented here provide a generative framework to address a fundamental paradox in contemporary democratic politics: how modern elections can experience sudden, late-stage reversals that diverge from eve-of-election survey forecasts despite the documented stability of voters' foundational political convictions. By developing an empirical digital twin of the 2026 Brazilian presidential contest, the model accounts for how late-campaign volatility can emerge not from the wholesale restructuring of underlying belief systems, but from an acute dynamical decoupling between fast-moving expressed opinions $\mathbf{x}_i(t)$ and slow-moving ideological convictions $\mathbf{s}_i(t)$.

This formulation provides a mechanistic explanation for the two-round polling anomaly observed in 2026 Brazil. While pre-election survey discrepancies in one-round races typically stem from a single late-breaking mobilization, the occurrence of sequential late-stage swings in opposing directions across sequential rounds reflects the dynamic properties of two-timescale opinion systems. Traditional polling methodologies sample top-of-the-head public opinions $\mathbf{x}_i(t)$ at discrete intervals under an implicit assumption of quasi-static equilibrium. In contemporary communication architectures, however, dark-social cascades accelerate during the final 48 hours of campaigning, precisely when statutory blackout periods and methodological field cutoffs suspend polling updates. Bridging Converse's belief constraint paradigm with Zaller's Receive-Accept-Sample model, the two-timescale formulation couples opinions to convictions through an elastic hyperbolic-tangent tether. Transient campaign blitzes displace expressed opinions $\mathbf{x}_i$ away from internal baselines, while slow-crystallizing convictions $\mathbf{s}_i$ exert a restorative drag that pulls opinions back once informational shocks subside. Consequently, an acute campaign blitz can produce sharp short-term voting swings that alter ballot outcomes without requiring deep ideological conversion, explaining how pre-election surveys can miss late-breaking inflection points.

My results suggest that this late-campaign electoral vulnerability is fundamentally governed by cognitive negativity asymmetry. In the calibrated model, candidate aversion accumulates with rapid sensitivity while decaying with extreme hysteresis ($\lambda_R = 0.14, \delta_R = 0.04$), producing a kinetic acquisition-to-decay ratio of $\lambda_R / \delta_R = 3.5\times$. In contrast, positive candidate admiration accumulates gradually and decays rapidly ($\lambda_A = 0.12, \delta_A = 0.06$). When combined with voting decision utilities that penalize candidate aversion more heavily than they reward positive affinity ($\theta_{\mathrm{rej}} = 1.65$), this kinetic asymmetry confers an inherent structural advantage upon negative disinformation blitzes. Positive candidate enthusiasm requires weeks of sustained communicative cultivation; in contrast, visceral candidate aversion can be synthesized within a 48-hour window, creating psychological scars that linger well beyond election day.

Furthermore, the multi-layer architecture exposes a structural pathology in contemporary platform curation. Commercial recommendation algorithms optimize for gross user engagement by weighting affirmative reactions and hostile outrage symmetrically, rewarding moral-emotional controversy along the diagonal where partisan admiration meets out-group outrage. Under this curation logic, conventional democratic defenses systematically misfire. Institutional newsrooms and regulatory courts operate under a wide separation of operational timescales ($\tau_{\mathrm{turnaround}} \gg \tau_{\mathrm{social}}$). In practice, defensive counter-speech, outrage debunking, and adversarial quote-tweets supply the out-group friction signals that recommendation engines use to distribute opponent messages into unaligned discovery feeds.

The second-round turnaround suggests that democratic resilience in algorithmic environments benefits from active socio-technical engagement rather than passive legalism alone. By translating the operational doctrine of \emph{Janonismo Cultural} into general graph-theoretic mechanisms, the model illustrates how coordinated socio-technical counter-strategies can systematically neutralize engagement-driven offensives. An outrage-starvation embargo (T1) deprives recommendation engines of engagement telemetry, collapsing opponent reach by $48\%$. This observation resonates with recent agent-based analyses of strategic behavioral norms in peer-to-peer networks, which indicate that simple shifts in local interaction rules and cooperative or defensive norms, such as withholding reactive engagement rather than amplifying friction, can nonlinearly transform macroscopic collective outcomes \citep{alvarenga2026power}. In parallel, monothematic bandwidth concentration (T2) exploits superlinear returns in platform discovery functions, quadrupling viral reach over dispersed messaging. Loss-aversion framing (T3) activates Kahneman--Tversky loss weighting (a 2.25-fold aversion to losses relative to equivalent gains), mobilizing cross-pressured working-class swing voters. Finally, molecular encrypted cascades (T4), physical turnout locking (T5), and attention-ceiling saturation (T6) bypass open-feed algorithms, insulate voters from late-stage blitzes, and cross the majority threshold.

Beyond domestic electoral competition, these platform dynamics carry critical geopolitical implications for digital sovereignty and transnational information flows. When the primary channels of democratic deliberation rely on proprietary recommendation systems controlled by foreign private technology corporations outside domestic regulatory jurisdiction, electoral outcomes become structurally vulnerable to opaque curation parameters and corporate content moderation choices. During recent Brazilian electoral cycles, major platforms such as Meta faced public scrutiny, parliamentary inquiries, and regulatory scrutiny by electoral courts regarding asymmetric advertising reach and perceived partisan disparities in algorithmic amplification. Moreover, because the scale-free topology of global social platforms aggregates engagement signals across borders, large-scale, synchronized domestic networks can generate cross-national algorithmic spillovers. This dynamic was empirically demonstrated when the 10-million-member \emph{Ex\'ercito do Lula} messaging network pivoted its coordinated mobilization infrastructure toward amplifying United States Democratic campaign narratives. Because recommendation engines weigh high-velocity interaction cascades without intrinsic regard for geographic origin, such massive coordinated domestic activity can alter global relevance scores ($R_m$) and cross algorithmic visibility thresholds in other countries, depending on how platform classifiers partition political content across linguistic and national boundaries.

Several modeling boundaries should be highlighted. First, the agent-based simulations treat interpersonal multiplex ties and platform follower networks as stationary topologies over the 50-day campaign horizon. While empirical network rewiring is negligible over weeks, multi-year electoral polarization inevitably induces selective tie severance and echo-chamber hardening. Second, the framework models political attitudes within a continuous metric concept hypercube $\Omega = [-1, 1]^4$. Although this captures core policy dimensions with high quantitative fidelity, it abstracts away linguistic nuances, narrative irony, and generative AI synthetic media. Future work should integrate large language model (LLM) agents capable of generating and interpreting unconstrained multimodal political rhetoric. Third, the second-round counterfactuals assume bounded counter-adaptation by political opponents, capturing short-term tactical surprises. In iterative electoral cycles, rival campaigns will adapt their own playbooks, sparking dynamic co-evolutionary arms races.

Notwithstanding these boundaries, the core behavioral mechanisms identified in this study, the two-timescale opinion-conviction tether, cognitive negativity hysteresis, and symmetric outrage curation, represent generalizable features of contemporary information ecosystems. Addressing democratic vulnerability requires structural platform reforms, such as deprecating out-group hostility signals in recommendation rankings, paired with the strategic capacity of democratic campaigns to defend popular majorities within algorithmic environments.

\section*{Methods}
\label{sec:methods}

\subsection*{Multidimensional concept space}
Citizens, political candidates, and media outlets are positioned within a $D$-dimensional continuous concept hypercube $\Omega = [-1, 1]^D$ ($D = 4$), whose coordinate axes represent the primary cleavages structuring contemporary democratic deliberation in Brazil: state economic redistribution and fiscal policy ($d=0$), moral, cultural, and religious customs ($d=1$), environmental governance and Amazonian preservation ($d=2$), and public security and firearms regulation ($d=3$) (Supplementary Note~1). The normalized distance between any two belief vectors $\mathbf{u}, \mathbf{v} \in \Omega$ is normalized by the maximal Euclidean diameter $2\sqrt{D}$:
\begin{equation}
    d(\mathbf{u}, \mathbf{v}) = \frac{1}{\sqrt{D}} \|\mathbf{u} - \mathbf{v}\|_2 \in [0, 2],
    \label{eq:methods_distance}
\end{equation}
such that individual conviction extremity is measured by the metric distance from the ideological origin:
\begin{equation}
    E(\mathbf{u}) \equiv d(\mathbf{u}, \mathbf{0}) = \frac{1}{\sqrt{D}} \|\mathbf{u}\|_2 \in [0, 1].
    \label{eq:methods_extremity}
\end{equation}
Ideological cross-issue constraint is enforced via a row-normalized cognitive coupling matrix $\mathbf{C} = [C_{dd'}] \in \mathbb{R}^{D \times D}$ satisfying decaying off-diagonal constraints (Supplementary Note~1). The strong cross-coupling between moral conservatism ($d=1$), environmental deregulation ($d=2$), and punitive security ($d=3$) directly reflects empirical roll-call voting alignments in the Brazilian National Congress.

\subsection*{Urban multiplex network topology}
Following urban digital twin methodologies established in computational epidemiology \citep{lopes2022measuring}, everyday interpersonal interactions unfold across a five-layer multiplex network $\mathcal{M} = \{\mathcal{L}_1, \dots, \mathcal{L}_5\}$ stored in Compressed Sparse Row (CSR) format, capturing the structural domains of urban social life (Supplementary Note~1):
Layer 1 ($\mathcal{L}_1$, Households): Disjoint fully connected cliques calibrated to national census statistics ($2.79$ residents per household, strong tie weight $w_1 = 0.36$);
Layer 2 ($\mathcal{L}_2$, Workplaces and Schools): Circular regular lattices with localized peer clustering ($\langle k^{[2]} \rangle = 10$, tie weight $w_2 = 0.26$);
Layer 3 ($\mathcal{L}_3$, Religious Congregations and Community Groups): Bipartite affiliation networks linking citizens to neighborhood churches and civic associations ($\langle k^{[3]} \rangle = 8$, weight $w_3 = 0.14$);
Layer 4 ($\mathcal{L}_4$, Friendship Circles): Small-world graphs with Watts--Strogatz rewiring probability $p_{\mathrm{ws}} = 0.15$ ($\langle k^{[4]} \rangle = 6$, weight $w_4 = 0.14$);
Layer 5 ($\mathcal{L}_5$, Public Transit and Incidental Weak Ties): Erd\H{o}s--R\'enyi random graphs representing transient weak ties in public transit ($\langle k^{[5]} \rangle = 4$, weight $w_5 = 0.10$).

Composite inter-agent edge weights $W_{ij} = \sum_{\ell=1}^5 w_\ell A_{ij}^{[\ell]} / \sum_k \sum_\ell w_\ell A_{ik}^{[\ell]}$ yield an empirical mean interpersonal degree of $\langle k_{\mathrm{comp}} \rangle \approx 30.8$ contacts per citizen. Inter-cohort contact is governed by the empirical contact matrix $M_{ab}$ across six adult age brackets, balancing generational clustering in workplaces against cross-generational mixing within households and religious assemblies (Supplementary Note~1).

\subsection*{Two-timescale attitude dynamics}
At each discrete integration step $\Delta t = 0.15$, citizen $i$ engages in peer deliberation with neighbors $j \in \mathcal{N}_i$ satisfying the bounded-confidence acceptance condition $d_{ij}(t) = \frac{1}{\sqrt{D}}\|\mathbf{x}_j(t) - \mathbf{x}_i(t)\|_2 \le \epsilon_i + 0.5\,u_i(t)$. Effective peer weights combine Relative Agreement (RA) uncertainty overlap $\phi_{ij}^{\mathrm{RA}}(t)$ and Biased Assimilation (BA) cosine similarity $\phi_{ij}^{\mathrm{BA}}(t)$:
\begin{equation}
    \tilde{w}_{ij}(t) = W_{ij} \big(0.40 + 0.60\,\phi_{ij}^{\mathrm{RA}}(t)\big) \phi_{ij}^{\mathrm{BA}}(t) \, \mathbb{I}\!\big(d_{ij}(t) \le \epsilon_i + 0.5\,u_i(t)\big).
    \label{eq:methods_peer_weight}
\end{equation}
The interpersonal social shift $\Delta \mathbf{x}_i^{\mathrm{soc}}(t) = \sum_{j} \tilde{w}_{ij}(\mathbf{x}_j - \mathbf{x}_i) / \sum_j \tilde{w}_{ij}$ updates current opinions through the Friedkin--Johnsen formulation:
\begin{equation}
    \mathbf{x}_i^{\mathrm{int}}(t) = \mathrm{clip}_{[-1, 1]^D}\!\Big(\mathbf{x}_i(t) + \Delta t \big[(1 - \alpha_i)(\mathbf{s}_i(t) - \mathbf{x}_i(t)) + \alpha_i \mathbf{C}\Delta \mathbf{x}_i^{\mathrm{soc}}(t)\big] + \boldsymbol{\xi}_i(t)\Big),
    \label{eq:methods_fj}
\end{equation}
with intrinsic cognitive noise $\boldsymbol{\xi}_i(t) \sim \mathcal{N}(\mathbf{0}, \sigma_{\xi}^2 \mathbf{I}_D)$ ($\sigma_{\xi} = 0.012$). When an exogenous factual impulse $\mathbf{b}_{\mathrm{fact}}(t)$ is present, opinions pull toward it with sensitivity $\gamma_i(t) = \gamma_0 (1.30 - 0.72 E_i)(1.18 - 0.48\tilde{a}_i)(0.85 + 0.30\,\mathrm{SES}_i)$, yielding intermediate opinion $\tilde{\mathbf{x}}_i(t)$.

Opinions cannot depart indefinitely from internal convictions. The maximal allowable departure is defined by an individual tether radius $\Delta_{\max, i}(t) = \Delta_0 (1.25 - 0.55\tilde{a}_i)(1.30 - 0.75 E_i)(0.90 + 0.25\,\mathrm{SES}_i)$. Radial elastic restraint is enforced via hyperbolic-tangent bounding:
\begin{align}
    g_i^{\mathrm{eff}}(t) &= \Delta_{\max, i}(t) \tanh\!\left(\frac{\frac{1}{\sqrt{D}}\|\tilde{\mathbf{x}}_i(t) - \mathbf{s}_i(t)\|_2}{\Delta_{\max, i}(t)}\right), \label{eq:methods_tether_radius} \\
    \mathbf{x}_i(t+1) &= \mathrm{clip}_{[-1, 1]^D}\!\left(\mathbf{s}_i(t) + \frac{g_i^{\mathrm{eff}}(t)}{g_i(t) + 10^{-12}}\big(\tilde{\mathbf{x}}_i(t) - \mathbf{s}_i(t)\big)\right). \label{eq:methods_tether_update}
\end{align}
Convictions adapt over substantially slower timescales through baseline crystallization drift $\eta_{\mathrm{cryst}, i}(t)$ and sustained-tension restructuring $\eta_{\mathrm{sust}, i}(t)$:
\begin{align}
    \eta_{\mathrm{cryst}, i}(t) &= \eta_0 [1 + (t/\tau_c)^{1.15}]^{-1}(1.25 - 0.75\tilde{a}_i)(1.22 - 0.68 E_i)(1.10 - 0.25\,\mathrm{SES}_i), \label{eq:methods_cryst} \\
    \eta_{\mathrm{sust}, i}(t) &= 0.28\eta_0 (g_i^{\mathrm{eff}} / \Delta_{\max, i})^2(1.18 - 0.58\tilde{a}_i)(1.20 - 0.62 E_i), \label{eq:methods_sust}
\end{align}
with baseline drift $\eta_0 = 0.0024 \text{ step}^{-1}$ and crystallization timescale $\tau_c = 280$ steps:
\begin{equation}
    \mathbf{s}_i(t+1) = \mathrm{clip}_{[-1, 1]^D}\!\Big(\mathbf{s}_i(t) + \big[\eta_{\mathrm{cryst}, i}(t) + \eta_{\mathrm{sust}, i}(t)\big]\big(\mathbf{x}_i(t+1) - \mathbf{s}_i(t)\big)\Big).
    \label{eq:methods_conviction}
\end{equation}

\subsection*{Mass media and digital platform rules}
Media conglomerates $g \in \{1, \dots, M\}$ operate across five modalities $v \in \{0, 1, 2, 3, 4\}$ (Broadcast Television, Radio, Print, Outdoor Signage, and Rallies) with demographic reach probabilities $P_v(a_i, \mathrm{SES}_i)$ (Supplementary Note~2). When a factual event $\mathbf{f}_k(t) \in \Omega$ occurs, outlet $g$ projects the notice toward its ideological center $\mathbf{m}_g \in \Omega$ through tendentiousness parameter $\tau_g \in [0, 1]$:
\begin{equation}
    \mathbf{y}_{k,g,v}(t) = (1 - \eta_{g,v})\,\mathbf{f}_k(t) + \eta_{g,v}\,\mathbf{m}_g + \boldsymbol{\xi}_{k,g,v}, \qquad \eta_{g,v} = \mathrm{clip}_{[0, 1]}(\tau_g \mu_v),
    \label{eq:methods_media_projection}
\end{equation}
where $\mu_v$ is modal baseline reach and $\sigma_{\mathrm{media}} = 0.03$. Published notices decay exponentially as $I_v(t+1) = I_v(t)\exp(-\delta_v)$. Citizens select active notices via distance-weighted softmax selective exposure:
\begin{equation}
    P(g, v \mid i) = \frac{P_v(a_i, \mathrm{SES}_i) \cdot I_{g,v}(t) \cdot \exp\!\left(-\beta_{\mathrm{media}} d(\mathbf{x}_i(t), \mathbf{y}_{g,v}(t))\right)}{\sum_{g',v'} P_{v'}(a_i, \mathrm{SES}_i) \cdot I_{g',v'}(t) \cdot \exp\!\left(-\beta_{\mathrm{media}} d(\mathbf{x}_i(t), \mathbf{y}_{g',v'}(t))\right) + \kappa_0},
    \label{eq:methods_selective_exposure}
\end{equation}
updating opinions within credibility-expanded confidence bounds $d(\mathbf{x}_i, \mathbf{y}_v) \le \epsilon_i (1 + 1.10\,\gamma_v)$ with coupling efficiency $\kappa_{\mathrm{trad}} = 0.08$ (Supplementary Note~2).

Digital media encompasses four platform modalities $p \in \{0, 1, 2, 3\}$ (Instagram/TikTok, YouTube/X, Facebook, and WhatsApp). Content creators accumulate followers via preferential attachment $\Pi_k \propto k_i (0.65 + 0.70\tau_{\mathrm{inf}, k})$, generating scale-free in-degrees $P(k) \sim k^{-1.15}$ with an empirical Gini coefficient of $0.73$ (Supplementary Note~3). Feed recommendation engines score candidate posts via homophily weighting $q_{ij} \propto (d(\mathbf{x}_i, \mathbf{y}_j) + \delta)^{-\beta_{\mathrm{cur}}}(1 + 1.5 e_j)\ln(1 + \mathrm{pop}_j)(0.5 + 0.5\ln(1 + k_{\mathrm{inf}}))$ with homophily exponent $\beta_{\mathrm{cur}} = 2.0$.

Reactions branch into affirmative likes ($d \le \epsilon_{\mathrm{like}, i}$, assimilating toward $\mathbf{y}_j$) and moral outrage ($d \ge \tau_{\mathrm{dislike}, i}$, triggering backfire repulsion away from $\mathbf{y}_j$). Algorithmic relevance scores compound dual-valence engagement symmetrically:
\begin{equation}
    R_m = \mathrm{Likes}_m + \omega \cdot \mathrm{Dislikes}_m, \qquad \omega = 1.0,
    \label{eq:methods_relevance}
\end{equation}
maximizing viral discovery along the diagonal of controversy (Supplementary Note~3). Information propagates through a three-stage cascade: open-feed broadcast injection, algorithmic recommendation discovery, and dark-social closed micro-group forwarding ($K \le 256$ participants per group, peer trust multiplier $\tau_{\text{WA}} = 0.85$, $2.8\times$ persuasion efficiency).

\subsection*{Electoral decision calculus and calibration}
Candidates $k \in \{L, F, O\}$ occupy fixed ideological coordinates $\mathbf{c}_k \in \Omega$. Candidate affinity $A_{ik}(t)$ and aversion $R_{ik}(t)$ evolve with asymmetric parameters ($\lambda_A = 0.12, \delta_A = 0.06; \lambda_R = 0.14, \delta_R = 0.04$):
\begin{align}
    A_{ik}(t+1) &= (1 - \delta_A)\,A_{ik}(t) + \lambda_A \max\!\left(0, \frac{0.62 - d_{ik}(t)}{0.62}\right)^2, \label{eq:methods_adm} \\
    R_{ik}(t+1) &= (1 - \delta_R)\,R_{ik}(t) + \lambda_R \max\!\left(0, \frac{d_{ik}(t) - 0.78}{1.22}\right)^{1.35}. \label{eq:methods_rej}
\end{align}
Net candidate utility is penalized quadratically by aversion:
\begin{equation}
    U_{ik}(t) = A_{ik}(t) - \theta_{\mathrm{rej}} R_{ik}(t)^2, \qquad \theta_{\mathrm{rej}} = 1.65.
    \label{eq:methods_utility}
\end{equation}
Conditional on voting, citizen $i$ chooses candidate $k \in \{L, F, O\}$ via multinomial logit:
\begin{equation}
    P(\text{Vote } k \mid i, \text{turnout}) = \frac{\exp(\beta_{\mathrm{vote}} U_{ik}(t))}{\sum_{m \in \{L, F, O\}} \exp(\beta_{\mathrm{vote}} U_{im}(t))},
    \label{eq:methods_logit}
\end{equation}
with choice sensitivity $\beta_{\mathrm{vote}} = 3.2$. Turnout probability is modeled via a logistic threshold:
\begin{equation}
    P(\text{Turnout} \mid i) = \frac{1}{1 + \exp\!\big(-\mu_{\mathrm{turn}}\big[\max_k U_{ik}(t) - (\tau_{\mathrm{abst}} - 0.22\,\mathrm{SES}_i + \delta_{\mathrm{transit}})\big]\big)},
    \label{eq:methods_turnout}
\end{equation}
with steepness $\mu_{\mathrm{turn}} = 4.5$ and free transit subsidy $\delta_{\mathrm{transit}} = -0.15$ (Supplementary Note~4).

To maintain rigorous methodological transparency, all model parameters are explicitly partitioned into three distinct classes (Supplementary Table~1 and Supplementary Table~2):

First, \emph{parameters fixed a priori from Brazilian demographic, institutional, and media records} are calibrated directly from empirical benchmark sources without free optimization:
(1) Citizen age distribution $a_i \in [16, 90]$ sampled from the official IBGE 2022 Census pyramid with Gaussian smoothing ($\sigma_a = 1.35$~years), socioeconomic status $\mathrm{SES}_i \sim \mathrm{Beta}(2.0, 3.0)$, and average household size of $2.79$ residents per domicile in Layer $\mathcal{L}_1$ \citep{ibge2022censo};
(2) Workplace and educational cohort sizes ($\langle k^{[2]} \rangle = 10$) drawn from INEP 2023 School Census and MTE/RAIS 2023 Formal Labor Registries \citep{inep2023censo, mte2023rais};
(3) Traditional media reach curves for television, radio, print, and outdoor signage calibrated to Kantar IBOPE Media Book 2023 audience surveys \citep{kantar2023inside};
(4) Digital platform penetration profiles across Instagram/TikTok, YouTube/X, Facebook, and WhatsApp ($K \le 256$ forwarding limit) derived from the Reuters Institute Digital News Report 2023 \citep{reuters2023digital};
(5) Statutory electoral parameters established by the Superior Electoral Court (TSE) under Lei n\textordmasculine~9.504/1997, comprising daily HGPE broadcast minutes (Lula 5m~31s vs.\ Fl\'avio 4m~20s in Round~1; strict 5m~00s parity in Round~2), presidential spending ceilings ($\mathrm{R}\$$~88.9M in Round~1; $\mathrm{R}\$$~133.4M in Round~2), and party allocations from the $\mathrm{R}\$$~4.96B FEFC public pool (Supplementary Note~4);
(6) Candidate ideological centroids in concept space $\Omega$: Lula $\mathbf{c}_L = (-0.25, -0.20, -0.35, -0.15)^{\top}$, Fl\'avio Bolsonaro $\mathbf{c}_F = (+0.65, +0.70, +0.60, +0.65)^{\top}$, and Third-Way Centrists $\mathbf{c}_O = (+0.20, +0.25, +0.10, +0.30)^{\top}$, derived from party manifesto content analyses and legislative roll-call matrices.

Second, \emph{parameters anchored in established cognitive and complex network literature} are set according to documented behavioral benchmarks:
(1) Friedkin--Johnsen social influence weighting and row-stochastic cognitive issue coupling matrix $\mathbf{C}$ \citep{friedkin1990social, friedkin2016network};
(2) Deffuant--Weisbuch bounded confidence threshold $\bar{\epsilon} = 0.45$ and Relative Agreement uncertainty dynamics \citep{deffuant2000mixing, deffuant2002how};
(3) Lord--Ross--Lepper biased assimilation cosine weighting $b_i$ \citep{lord1979biased, dandekar2013biased};
(4) Kahneman--Tversky prospect theory loss aversion multiplier $\lambda_{\mathrm{loss}} \in [2.0, 2.5]$ \citep{kahneman1979prospect};
(5) Watts--Strogatz small-world rewiring probability $p_{\mathrm{ws}} = 0.15$ in friendship networks ($\mathcal{L}_4$) \citep{watts1998collective};
(6) Granovetter weak-tie mixing in transit layer $\mathcal{L}_5$ \citep{granovetter1973strength};
(7) Barab\'asi--Albert preferential attachment scaling exponent $\gamma_{\mathrm{pl}} = 1.15$ and Gini coefficient $0.73$ for digital influencers \citep{barabasi1999emergence, bianconi2001competition};
(8) Hyperbolic-tangent elastic tethering $g_i^{\mathrm{eff}}$ enforcing Sherif--Hovland latitudes of rejection around convictions \citep{sherif1961social, howe2017attitude}.

Third, \emph{free parameters fitted to longitudinal 2026 Datafolha tracking surveys and TSE ballot totals} are estimated via bounded grid search and Markov Chain Monte Carlo (MCMC) sampling:
(1) Fast opinion discrete integration step size $\Delta t = 0.15$ and cognitive perturbation noise $\sigma_{\xi} = 0.012$;
(2) Baseline conviction adaptation drift rate $\eta_0 = 0.0024 \text{ step}^{-1}$ and crystallization timescale $\tau_c = 280$ simulation steps;
(3) Direct exogenous fact pull coefficient $\gamma_0 = 0.055 \text{ step}^{-1}$;
(4) Candidate affinity and aversion kinetic rates ($\lambda_A = 0.12, \delta_A = 0.06; \lambda_R = 0.14, \delta_R = 0.04$), establishing the key kinetic asymmetry ratio $\lambda_R / \delta_R = 3.5\times$ and aversion penalty $\theta_{\mathrm{rej}} = 1.65$;
(5) Multinomial logit voting sensitivity $\beta_{\mathrm{vote}} = 3.2$, base abstention threshold $\tau_{\mathrm{abst}} = 0.38$, logistic steepness $\mu_{\mathrm{turn}} = 4.5$, and transit subsidy $\delta_{\mathrm{transit}} = -0.15$.

The optimization objective minimizes weighted squared error across all survey waves and official election totals:
\begin{equation}
    L(\boldsymbol{\theta}) = \sum_{w=1}^W \sum_{k \in \{L, F, O\}} \big(V_k^{\mathrm{model}}(t_w; \boldsymbol{\theta}) - V_k^{\mathrm{Datafolha}}(t_w)\big)^2 + \omega_{\mathrm{ballot}} \sum_{k} \big(V_k^{\mathrm{model}}(t_{\mathrm{election}}; \boldsymbol{\theta}) - V_k^{\mathrm{TSE}}\big)^2,
    \label{eq:methods_loss_function}
\end{equation}
with ballot weight $\omega_{\mathrm{ballot}} = 3.0$, yielding $\mathrm{RMSE} = 0.85$ percentage points ($R^2 = 0.9955$) across Round~1 and $\mathrm{RMSE} = 0.56$ percentage points ($R^2 = 0.9128$) across Round~2. Profile likelihood intervals, Hessian conditioning, parameter identifiability, and Sobol global sensitivity analyses are detailed extensively in Supplementary Note~4 and Supplementary Note~5. Operational translations of the Janones tactical suite into graph mechanics are reported in Supplementary Table~3.

\subsection*{Computational implementation}
Simulations were executed across $S = 100$ independent Monte Carlo realizations at population scale $N = 100{,}000$ agents using just-in-time (JIT) compiled Numba 0.59 kernels in Python 3.11 with parallel multi-threading across 64 CPU cores (AMD EPYC 7763, 2.45 GHz, 512 GB RAM) (Supplementary Note~5). All agent attributes are laid out as contiguous 1D and 2D C-ordered NumPy arrays in CPU cache memory, with network topology represented in CSR format. At $N = 100{,}000$, a 100-step trajectory executes in $4.2$ seconds. Municipal scaling experiments ($N = 10^4$ to $10^7$) confirm that finite-size stochasticity vanishes into a scale-invariant thermodynamic limit for $N \ge 10^5$ (Supplementary Note~5).

Macroscopic system trajectories are monitored through four statistical evaluation metrics:
(1) Multidimensional Esteban--Ray polarization index $P_{\mathrm{ER}}(t) = \kappa_{\mathrm{ER}} \sum_{c, c'} \pi_c^{1 + \alpha_{\mathrm{ER}}} \pi_{c'} d(\bar{\mathbf{x}}_c, \bar{\mathbf{x}}_{c'})$ ($\alpha_{\mathrm{ER}} = 1.30$, $\kappa_{\mathrm{ER}} = 1.0$) \citep{esteban1994measurement};
(2) Sarle's Bimodality Coefficient $\mathrm{BC} = (\gamma^2 + 1) / (\kappa + 3 \frac{(n-1)^2}{(n-2)(n-3)})$;
(3) Participation Ratio of the belief covariance matrix $\mathrm{PR}(\Sigma) = (\mathrm{Tr}\,\Sigma)^2 / \mathrm{Tr}(\Sigma^2) \in [1, D]$ measuring effective ideological dimensionality \citep{roy1953useful};
(4) Mean absolute pairwise issue alignment $\langle |\rho| \rangle = \frac{2}{D(D-1)} \sum_{k < l} |\rho_{kl}|$. All time-series trajectories report ensemble means with $95\%$ confidence intervals ($\bar{\mu} \pm 1.96\,\hat{\sigma}/\sqrt{S}$) or full simulation ranges $[\min, \max]$.

\section*{Acknowledgements}
I thank Luana Alvarenga, Renata Mota, and Fernando Pureza for insightful discussions and feedback. I acknowledge the high-performance computing resources provided by the N{\'u}cleo de Processamento de Alto Desempenho (NPAD) at the Federal University of Rio Grande do Norte (UFRN). This work received financial and institutional support from the Coordena{\c{c}}{\~a}o de Aperfei{\c{c}}oamento de Pessoal de N{\'i}vel Superior (CAPES), the Conselho Nacional de Desenvolvimento Cient{\'i}fico e Tecnol{\'o}gico (CNPq), and the Pr{\'o}-Reitoria de Pesquisa of the Federal University of Rio Grande do Norte (UFRN/PROPESQ).

\section*{Data availability}
The empirical demographic, polling, broadcast airtime, and election result datasets supporting the findings of this study are archived and publicly available at \url{https://github.com/cesarrennocosta/janones}.

\section*{Code availability}
All simulation code, Numba JIT kernels, graph generation scripts, and figure plotting pipelines are open source under the MIT License at \url{https://github.com/cesarrennocosta/janones}.

\bibliographystyle{naturemag}
\bibliography{references}

\end{document}
